% GENERATED FILE. Edit canonical lessons, then run npm run generate:books.
% canonical-source-hash: fnv1a64:128547cb2e751c5d
% canonical-lessons: FR-C463-verglas, FR-C463-givre, FR-C463-geler, FR-C463-degeler, FR-C463-grelotter, FR-R463-revision-456-463

\chapter{Le grand froid}
\label{ch:fr-le-grand-froid}
\hlchaptermodality{463}

\begin{chapteropening}
\textbf{By the end of this chapter:} \emph{I can talk about freezing weather: black ice, frost, freezing, thawing and shivering.}

Complete the last lesson of chapter 463: I can talk about freezing weather: black ice, frost, freezing, thawing and shivering.
\end{chapteropening}

\section[le verglas]{le verglas — the black ice}

\label{lesson:FR-C463-verglas}



\begin{quote}
Before the new one: say the French for to rumble, then the French for the gust.
\end{quote}

\subsection*{You'll want to know: le verglas}
\textbf{le verglas} — \textquotedblleft{}the black ice\textquotedblright{}.

\begin{quote}
Avec le verglas, les voitures roulent très lentement.
\end{quote}

\begin{quote}
\emph{Because of the black ice, cars are going very slowly.}
\end{quote}

\emph{Une route verglacée} — an icy road.

\begin{cousinweb}
\emph{verglas} is built from \emph{verre} 'glass', from Latin \emph{vitrum}, and \emph{glace} 'ice', from Latin \emph{glaciēs} (EtymDB): ice as smooth as glass.
\end{cousinweb}

\begin{grammarlens}[title={What you've built}]
One more word for telling, explaining and describing at B1.
\end{grammarlens}

\subsection*{Your turn}
\begin{itemize}
\raggedright
  \item \emph{Say it:} \emph{le verglas}
  \item \emph{Say it:} \emph{le verglas}, once more
  \item \emph{Say it:} a sentence of your own with \emph{le verglas}
  \item \emph{Recall:} say the French for to rumble, then the French for the gust, then say \emph{le verglas} again
\end{itemize}

\begin{tcolorbox}[breakable,colback=teal!4,colframe=teal!35!black,title={Before you move on}]
What does le verglas mean? (\textquotedblleft{}The black ice\textquotedblright{}.) Say it once more.
\end{tcolorbox}
\section[le givre]{le givre — the frost (on surfaces), the hoarfrost}

\label{lesson:FR-C463-givre}



\begin{quote}
Before the new one: say the French for the gust, then the French for the black ice.
\end{quote}

\subsection*{You'll want to know: le givre}
\textbf{le givre} — \textquotedblleft{}the frost (on surfaces), the hoarfrost\textquotedblright{}.

\begin{quote}
Le matin, il y a du givre sur le pare-brise.
\end{quote}

\begin{quote}
\emph{In the morning there is frost on the windscreen.}
\end{quote}

White frost on windows and grass.

\begin{cousinweb}
The origin of \emph{givre} is not settled; it may come from a language spoken in the region before Latin. EtymDB lists it without a source.
\end{cousinweb}

\begin{grammarlens}[title={What you've built}]
One more word for telling, explaining and describing at B1.
\end{grammarlens}

\subsection*{Your turn}
\begin{itemize}
\raggedright
  \item \emph{Say it:} \emph{le givre}
  \item \emph{Say it:} \emph{le givre}, once more
  \item \emph{Say it:} a sentence of your own with \emph{le givre}
  \item \emph{Recall:} say the French for the gust, then the French for the black ice, then say \emph{le givre} again
\end{itemize}

\begin{tcolorbox}[breakable,colback=teal!4,colframe=teal!35!black,title={Before you move on}]
What does le givre mean? (\textquotedblleft{}The frost (on surfaces), the hoarfrost\textquotedblright{}.) Say it once more.
\end{tcolorbox}
\section[geler]{geler — to freeze}

\label{lesson:FR-C463-geler}



\begin{quote}
Before the new one: say the French for the black ice, then the French for the frost.
\end{quote}

\subsection*{You'll want to know: geler}
\textbf{geler} — \textquotedblleft{}to freeze\textquotedblright{}.

\begin{quote}
Cette nuit, il a gelé et le lac est blanc.
\end{quote}

\begin{quote}
\emph{Last night it froze and the lake is white.}
\end{quote}

\emph{Il gèle} — it is freezing; \emph{je suis gelé} — I am frozen.

\begin{cousinweb}
\emph{geler} is inherited through Old French from Latin \emph{gelāre} 'to freeze', from \emph{gelū} 'frost' (EtymDB).
\end{cousinweb}

\begin{grammarlens}[title={What you've built}]
One more word for telling, explaining and describing at B1.
\end{grammarlens}

\subsection*{Your turn}
\begin{itemize}
\raggedright
  \item \emph{Say it:} \emph{geler}
  \item \emph{Say it:} \emph{geler}, once more
  \item \emph{Say it:} a sentence of your own with \emph{geler}
  \item \emph{Recall:} say the French for the black ice, then the French for the frost, then say \emph{geler} again
\end{itemize}

\begin{tcolorbox}[breakable,colback=teal!4,colframe=teal!35!black,title={Before you move on}]
What does geler mean? (\textquotedblleft{}To freeze\textquotedblright{}.) Say it once more.
\end{tcolorbox}
\section[dégeler]{dégeler — to thaw}

\label{lesson:FR-C463-degeler}



\begin{quote}
Before the new one: say the French for the frost, then the French for to freeze.
\end{quote}

\subsection*{You'll want to know: dégeler}
\textbf{dégeler} — \textquotedblleft{}to thaw\textquotedblright{}.

\begin{quote}
Au printemps, la rivière dégèle.
\end{quote}

\begin{quote}
\emph{In spring the river thaws.}
\end{quote}

\emph{Le dégel} — the thaw.

\begin{cousinweb}
\emph{dégeler} is built from \emph{dé-} and \emph{geler} 'to freeze' (EtymDB lists it), which is inherited from Latin \emph{gelāre}.
\end{cousinweb}

\begin{grammarlens}[title={What you've built}]
One more word for telling, explaining and describing at B1.
\end{grammarlens}

\subsection*{Your turn}
\begin{itemize}
\raggedright
  \item \emph{Say it:} \emph{dégeler}
  \item \emph{Say it:} \emph{dégeler}, once more
  \item \emph{Say it:} a sentence of your own with \emph{dégeler}
  \item \emph{Recall:} say the French for the frost, then the French for to freeze, then say \emph{dégeler} again
\end{itemize}

\begin{tcolorbox}[breakable,colback=teal!4,colframe=teal!35!black,title={Before you move on}]
What does dégeler mean? (\textquotedblleft{}To thaw\textquotedblright{}.) Say it once more.
\end{tcolorbox}
\section[grelotter]{grelotter — to shiver (with cold)}

\label{lesson:FR-C463-grelotter}



\begin{quote}
Before the new one: say the French for to freeze, then the French for to thaw.
\end{quote}

\subsection*{You'll want to know: grelotter}
\textbf{grelotter} — \textquotedblleft{}to shiver (with cold)\textquotedblright{}.

\begin{quote}
Sans manteau, le garçon grelotte à l'arrêt de bus.
\end{quote}

\begin{quote}
\emph{Without a coat the boy is shivering at the bus stop.}
\end{quote}

Only from cold.

\begin{cousinweb}
\emph{grelotter} is built on \emph{grelot} 'small round bell' (EtymDB), usually explained by an old phrase about teeth that rattle like a little bell.
\end{cousinweb}

\begin{grammarlens}[title={What you've built}]
One more word for telling, explaining and describing at B1.
\end{grammarlens}

\subsection*{Your turn}
\begin{itemize}
\raggedright
  \item \emph{Say it:} \emph{grelotter}
  \item \emph{Say it:} \emph{grelotter}, once more
  \item \emph{Say it:} a sentence of your own with \emph{grelotter}
  \item \emph{Recall:} say the French for to freeze, then the French for to thaw, then say \emph{grelotter} again
\end{itemize}

\begin{tcolorbox}[breakable,colback=teal!4,colframe=teal!35!black,title={Before you move on}]
What does grelotter mean? (\textquotedblleft{}To shiver (with cold)\textquotedblright{}.) Say it once more.
\end{tcolorbox}
\section[(dialogue)]{Retour sur les chapitres 456–463}

\label{lesson:FR-R463-revision-456-463}



\begin{quote}
Say the French for the stream and the pond.
\end{quote}

\subsection*{Your turn}
Say these aloud:

\begin{itemize}
\raggedright
  \item the stream, the pond, the marsh, to splash about, to stream
  \item the clearing, the bush, the fern, the moss, the bark
  \item the dune, the tide, the cove, the reef, the seashell
  \item the squirrel, the hedgehog, the wild boar, the stag, the doe
  \item the weather forecast, the forecast, cloudy, sunny, the thermometer
  \item the drizzle, to drizzle, the rain shower, the bright spell, to clear
  \item stormy, the lightning, the hail, to rumble, the gust
  \item the black ice, the frost, to freeze, to thaw, to shiver
\end{itemize}

\begin{tcolorbox}[breakable,colback=teal!4,colframe=teal!35!black,title={Before you move on}]
The stream? (\textbf{le ruisseau}.) The pond? (\textbf{l'étang}.) The marsh? (\textbf{le marais}.) To splash about? (\textbf{patauger}.) To stream? (\textbf{ruisseler}.) The clearing? (\textbf{la clairière}.) The bush? (\textbf{le buisson}.) The fern? (\textbf{la fougère}.) The moss? (\textbf{la mousse}.) The bark? (\textbf{l'écorce}.) The dune? (\textbf{la dune}.) The tide? (\textbf{la marée}.) The cove? (\textbf{la crique}.) The reef? (\textbf{le récif}.) The seashell? (\textbf{le coquillage}.) The squirrel? (\textbf{l'écureuil}.) The hedgehog? (\textbf{le hérisson}.) The wild boar? (\textbf{le sanglier}.) The stag? (\textbf{le cerf}.) The doe? (\textbf{la biche}.) The weather forecast? (\textbf{la météo}.) The forecast? (\textbf{les prévisions}.) Cloudy? (\textbf{nuageux}.) Sunny? (\textbf{ensoleillé}.) The thermometer? (\textbf{le thermomètre}.) The drizzle? (\textbf{la bruine}.) To drizzle? (\textbf{bruiner}.) The rain shower? (\textbf{l'averse}.) The bright spell? (\textbf{l'éclaircie}.) To clear? (\textbf{dégager}.) Stormy? (\textbf{orageux}.) The lightning? (\textbf{la foudre}.) The hail? (\textbf{la grêle}.) To rumble? (\textbf{gronder}.) The gust? (\textbf{la rafale}.) The black ice? (\textbf{le verglas}.) The frost? (\textbf{le givre}.) To freeze? (\textbf{geler}.) To thaw? (\textbf{dégeler}.) To shiver? (\textbf{grelotter}.)
\end{tcolorbox}
